\section{Catchphrases: когда помогают, а когда превращают персонажа в куклу}

Catchphrase помогает, когда она сжимает уже установленный характер в мгновенно узнаваемую форму (тот же принцип силуэта, 6.1, только для голоса), и вредит, когда становится заменой характера, а не его следствием.

"In the name of the Moon, I'll punish you!" Сейлор Мун (6.6) — рабочий пример именно потому, что фраза не существует отдельно от ритуала трансформации (1.3): она произносится в момент, когда обычная школьница буквально становится кем-то другим, и несёт вес этого перехода, а не просто узнаваемость сама по себе. Без трансформационного контекста та же фраза читалась бы как пустая бравада; внутри контекста — как заявление о готовности, подтверждённое действием, которое сейчас последует.

Более тонкая опасность — механическое повторение, уже названное в 2.6 применительно к ритуалам: фраза, которая звучит идентично при каждом появлении вне зависимости от ставок момента, быстро превращается из узнаваемости в тик, а персонаж — в куклу, дёргающую за одну и ту же ниточку. Тест простой: можно ли предсказать catchphrase персонажа, зная только его типаж, — если да, это ярлык, а не сигнатура, в отличие от фразы Сейлор Мун, которая невозможна вне конкретной истории конкретной девочки. Рабочее правило то же, что и для ритуалов: функция повторяется, слова — нет.
